\documentclass[11pt,letterpaper]{article}
\usepackage{technical-implementation}
% Shared non-semantic helper definitions; one semantic module is loaded by each root.
\newcommand{\PTSource}[1]{\par\smallskip{\footnotesize\color{Meta}\textbf{Book crosswalk:} #1}\par\smallskip}
\newcommand{\PTOutcome}[1]{\begin{tipbox}[Competency outcome]#1\end{tipbox}}
\newcommand{\PTEvidence}[1]{\begin{notebox}[Required evidence]#1\end{notebox}}
\newcount\ptlinecount
\newcommand{\PTLines}[1]{\par\ptlinecount=0\loop\ifnum\ptlinecount<#1\par\noindent\rule{\linewidth}{0.25pt}\vspace{0.22cm}\advance\ptlinecount by1\repeat\par}
\newcommand{\PTField}[1]{\noindent\textbf{#1}\quad\rule{0.66\linewidth}{0.3pt}\par\vspace{0.22cm}}
\newcommand{\PTRating}{\noindent\textbf{Rating:}\quad 0\quad 1\quad 2\quad 3\quad 4\qquad\textbf{Evidence ref:}\rule{0.30\linewidth}{0.3pt}\par\vspace{0.14cm}}
\newcommand{\PTTask}[2]{\par\medskip\noindent\textbf{Task #1.} #2\par}

\setDocTitle{Parsing and Compiler Engineering: 24-Week Curriculum}
\setDocSubtitle{Dependency ordered with evidence-based progression}
\setDocVersion{1.0}\setDocStatus{Engineering education / source-aligned}
\setDocOwner{Compiler Engineering}
\setDocDate{30 October 2026}
\setDocSummary{Prerequisite-based learning curriculum.}
\begin{document}\makedoctitle\makedocsummary
\tableofcontents\clearpage
\section{Entry competencies}
Comfort with recursion, stacks, queues, sets/maps, elementary discrete mathematics, big-O reasoning and basic testing in Python/C++. If gaps exist, complete the matching workbook before the relevant stage.
\section{Curriculum and dependencies}
\small\begin{longtable}{p{0.12\linewidth}p{0.23\linewidth}p{0.10\linewidth}p{0.23\linewidth}p{0.24\linewidth}}\toprule Weeks & Module & Book chapters & Core knowledge & Exit evidence\\\midrule\endhead
1--2 & Formal language foundations & 1, 2 & derivations, nullable, reachability, productivity & Complete grammar worksheets and identify empty language versus empty string.\\
3--4 & Parse structure and finite automata & 3, 5 & parse trees, ambiguity, NFA/DFA & Implement DFA recognizer and preserve accepted/rejected fixtures.\\
5--6 & General CF recognition & 4, 7 & CYK, chart invariants, Earley & Test CYK and Earley on the same finite sample corpus.\\
7--8 & Top-down parser construction & 6, 8 & left recursion, recursive descent, LL(1) & Produce FIRST/FOLLOW and a conflict-free predictive table.\\
9--10 & Bottom-up parser construction & 9 & LR items, closure/goto, conflicts & Produce LR state diagrams and a conflict report.\\
11--12 & Generalized and noncanonical methods & 10, 11 & left-corner, GLR/GLL and forest sharing & Explain ambiguity effects in a constrained prototype.\\
13--14 & Substring, intersection and parallel & 12, 13, 14 & suffix grammar, FSA intersection, work-depth & Prepare three comparative design analyses with limitations.\\
15--16 & Semantic extensions and recovery & 15, 16 & attributes, tree-based checks, diagnostics & Implement error reporting and a restricted semantic checker.\\
17--18 & Practical compiler architecture & 17 & AST, lexer/parser API, symbol/type analysis & Deliver an end-to-end mini-language compiler with source spans.\\
19--20 & Security hardening & 16, 17 & fuzzing, bounded resources, adversarial testing & Complete threat model, fuzz run and a resolved defect record.\\
21--22 & Reproducible benchmarking & 4, 7, 8, 9 & hypothesis, generator, sampling, reporting & Present raw samples and analyze empirical scaling.\\
23--24 & Capstone, literature and defense & 18 + review & source evaluation, trade-offs and technical writing & Defend parser selection with proof, tests, timings and a scored review.\\
\bottomrule\end{longtable}\normalsize
\section{Dependency and remediation policy}
A stage closes only when the learner produces accepted/rejected fixtures, a written invariant, and a peer-readable evidence artifact. Return to Ch. 2 when grammar transformations are unclear; Ch. 5 before automaton-heavy methods; Ch. 3 before comparing parsing directions. Advanced chapters are electives after the core deterministic and chart-based methods.
\section{Assessment gates}
\begin{enumerate}\item Gate A (week 4): formal grammar and DFA correctness.\item Gate B (week 10): CYK/Earley/LL/LR reasoning, including conflict analysis.\item Gate C (week 18): end-to-end compiler tests and error contract.\item Gate D (week 24): secure and reproducible capstone with measurements, security evidence and architecture review.\end{enumerate}

\end{document}
